% Source fingerprint: 89a41d06e175d9aa5ca92aa3ddd8667a41bbec1369852f31d74b225a69b2e5b8
% T09: raw TeX cells preserved; no numeric highlighting.
\begin{table}[htbp]
\centering\small\renewcommand{\arraystretch}{1.18}\setlength{\tabcolsep}{4pt}
\caption[Lipschitz、绝对连续与有界变差]{Lipschitz、绝对连续与有界变差。这里的一维变差按划分定义。第27章从分布导数测度重新组织该理论；函数代表与端点取值始终按正文约定。}
\label{b1:tab:visual-t09}
\begin{tabularx}{\linewidth}{@{}>{\raggedright\arraybackslash}p{3.1cm}>{\raggedright\arraybackslash}p{4.7cm}>{\raggedright\arraybackslash}X@{}}
\toprule
有限区间上的函数类 & 导数与积分恢复 & 不能逆推的例子 \\
\midrule
Lipschitz & 绝对连续；导数几乎处处存在且本性有界；积分恢复函数增量 & 绝对连续函数不必 Lipschitz，如平方根 \\
绝对连续 & 导数属于 \(L^1\)，且 \(f(t)-f(a)=\int_a^t f^{\prime}(s)\,\mathrm{d}s\) & 仅有连续性和几乎处处可微不足以保证积分恢复 \\
有界变差 & 划分增量绝对值之和有统一上界；导数可积，但可有奇异变化 & Cantor 函数连续、几乎处处导数为零，却不是常值 \\
一般连续函数 & 不保证有界变差或几乎处处可微 & 本章三角峰例子可使划分增量和发散 \\
\bottomrule
\end{tabularx}
\end{table}
